\documentclass[11pt]{article}
\usepackage[utf8]{inputenc}
\usepackage[T1]{fontenc}
\usepackage[french]{babel}
\usepackage{lmodern}
\usepackage{xurl}
\usepackage{hyperref}
\usepackage{enumitem}
\usepackage{geometry}
\geometry{margin=1in}

\title{Fiche de cours -- Informatique quantique}
\date{}

\begin{document}
\maketitle

Ce cours propose une introduction large et accessible à l’informatique quantique, avec des démonstrations pratiques utilisant l’écosystème \textbf{Qiskit} d’IBM. Il est conçu pour des participants issus d’horizons variés tels que des étudiants, développeurs et professionnels disposant de connaissances mathématiques de base (vecteurs, matrices) et en programmation Python. Nous apprendrons à implémenter des circuits quantiques élémentaires et à les exécuter sur de véritables ordinateurs quantiques mis à la disposition du grand public par IBM.

\section*{Content}

Ce cours abordera les thématiques suivantes:

\begin{itemize}[leftmargin=*,itemsep=0.4em]
\item Informatiques quantique et classique : avantages et motivations;
\item Notions essentielles de mécanique quantique appliquées à l’informatique : qubit, portes logiques quantiques, notation bra-ket, états superposés, mesure, \ldots
\item Circuits quantiques multi-qubits : intrication quantique, manipulation des portes logiques;
\item Programmation de circuits quantiques avec la librairie \textbf{Qiskit} : préparation d'un état intriqué et téléportation quantique;
\item Aperçu d'un algorithme quantique plus avancé avec \textbf{Qiskit} : algorithme de Grover;
\item Exécution de circuits sur un véritable ordinateur quantique;
\end{itemize}

En plus de ces parties, d'autres thématiques pourront être dispensées à la discrétion des participants:

\begin{itemize}[leftmargin=*,itemsep=0.4em]
\item Panorama des plateformes matérielles pour l’informatique quantique : technologies existantes et limitations actuelles;
\item CHSH et inégalités de Bell : interprétation de la superposition quantique;
\item Les nombres complexes pour l'informatique quantique;
\item Évolution des états quantiques et impossibilité de copier un état quantique;
\item Distribution quantique de clés de chiffrement.
\end{itemize}

\section*{Learning Outcomes}

À l’issue de cette formation, les participants auront :

\begin{itemize}[leftmargin=*,itemsep=0.4em]
\item Une compréhension élémentaire de l’informatique quantique;
\item La capacité d’implémenter des circuits quantiques simples avec Qiskit;
\item La capacité d’exécuter des algorithmes sur les ordinateurs quantiques d’IBM accessibles au public.
\end{itemize}

\section*{Training Method}

Cette formation alterne exposés théoriques et démonstrations pratiques.

\section*{Prerequisites}

\begin{itemize}[leftmargin=*,itemsep=0.4em]
\item Connaissances mathématiques de base des vecteurs et des matrices;
\item Programmation Python;
\end{itemize}

Lecture recommandée en algèbre linéaire (vecteurs, matrices):
\url{https://www.khanacademy.org/math/linear-algebra}

\end{document}
